\subsection{Schrödingergleichung für das Wasserstoffatom}
\label{subsec:schroedingergleichung_wasserstoffatom}

Der Wellencharakter des Elektrons\index{Wellencharakter!des Elektrons} führt
direkt zur Frage, wie sich ein Elektron\index{Elektron} mathematisch beschreiben
lässt. In der klassischen Mechanik wird die Bewegung eines Teilchens durch seine
Bahn beschrieben. Man gibt Ort und Geschwindigkeit an und berechnet daraus die
weitere Bewegung. In der Quantenmechanik\index{Quantenmechanik} ist diese
Vorstellung nicht mehr ausreichend.

Die zentrale Größe ist nicht mehr eine Bahnkurve, sondern die
Wellenfunktion\index{Wellenfunktion}. Sie wird meist mit \(\psi\) bezeichnet.
Die Wellenfunktion selbst ist nicht direkt messbar. Aus ihrem Betragsquadrat
ergibt sich jedoch die Wahrscheinlichkeit\index{Wahrscheinlichkeit}, ein
Elektron an einem bestimmten Ort zu finden:

\begin{equation}
	P(\vec r) \propto |\psi(\vec r)|^2 .
	\label{eq:wahrscheinlichkeit_psi_quadrat}
\end{equation}

Die Gleichung, mit der sich die Wellenfunktion eines nichtrelativistischen
Teilchens berechnen lässt, ist die Schrödingergleichung
\index{Schrödingergleichung}. Für stationäre Zustände besitzt sie die Form:

\begin{equation}
	\hat{H}\psi = E\psi .
	\label{eq:stationaere_schroedingergleichung}
\end{equation}

Dabei ist \(\hat{H}\)\index{Hamilton-Operator} der Hamilton-Operator. Er
beschreibt die Gesamtenergie des Systems. Die Größe \(E\)\index{Energie} ist
die Energie des jeweiligen Zustands. Die Wellenfunktion \(\psi\) beschreibt den
zugehörigen quantenmechanischen Zustand\index{Zustand!quantenmechanischer}.

Für das Wasserstoffatom\index{Wasserstoffatom} ist diese Gleichung besonders
wichtig. Das Wasserstoffatom besteht aus einem Proton\index{Proton} und einem
Elektron. Das Elektron befindet sich im elektrischen Feld des positiv geladenen
Atomkerns\index{Atomkern}. Zwischen Elektron und Proton wirkt die
Coulomb-Anziehung\index{Coulomb-Anziehung}.

In vereinfachter Form lautet die stationäre Schrödingergleichung für das
Elektron im Wasserstoffatom:
In dieser vereinfachten Darstellung wird der Atomkern als ruhend angenommen.
Die mathematischen Hintergründe dieser Gleichung werden in
Anhang~\ref{app:mathe_schroedingergleichung} und Anhang~\ref{subsec:mathematische_voraussetzungen_schroedinger}  genauer beschrieben.
\begin{equation}
	\left(
	-\frac{\hbar^2}{2m_\mathrm{e}}\nabla^2
	-\frac{e^2}{4\pi\varepsilon_0 r}
	\right)\psi
	=
	E\psi .
	\label{eq:schroedingergleichung_wasserstoff}
\end{equation}

Der erste Term

\begin{equation}
	-\frac{\hbar^2}{2m_\mathrm{e}}\nabla^2
	\label{eq:kinetischer_term_schroedinger}
\end{equation}

beschreibt die kinetische Energie\index{kinetische Energie} des Elektrons. Der
Operator \(\nabla^2\)\index{Laplace-Operator} ist der Laplace-Operator. Er
beschreibt, wie stark sich die Wellenfunktion räumlich verändert.

Der zweite Term

\begin{equation}
	-\frac{e^2}{4\pi\varepsilon_0 r}
	\label{eq:coulomb_potential_wasserstoff}
\end{equation}

beschreibt die potentielle Energie\index{potentielle Energie} des Elektrons im
Coulomb-Feld\index{Coulomb-Feld} des Protons. Das Minuszeichen zeigt, dass die
Kraft anziehend ist. Je kleiner der Abstand \(r\)\index{Abstand} zwischen
Elektron und Proton ist, desto stärker ist diese Wechselwirkung.

\begin{PhysicsBox}[Bedeutung der Schrödingergleichung]
	\label{box:bedeutung_schroedingergleichung}
	Die Schrödingergleichung ersetzt im Atom die klassische Bahn des Elektrons
	durch eine Wellenfunktion. Aus dieser Wellenfunktion ergeben sich die
	möglichen Energien und die Aufenthaltswahrscheinlichkeiten des Elektrons.
\end{PhysicsBox}

Die Lösungen der Schrödingergleichung für das Wasserstoffatom liefern nicht
beliebige Energien. Es sind nur bestimmte Energiewerte erlaubt. Man spricht von
diskreten Energieniveaus\index{Energieniveau!diskretes}. Damit erklärt die
Quantenmechanik, warum Atome nur bestimmte Lichtfrequenzen aufnehmen oder
aussenden können.

Gleichzeitig liefern die Lösungen räumliche Verteilungen für die
Aufenthaltswahrscheinlichkeit des Elektrons. Diese Verteilungen sind keine
klassischen Kreisbahnen. Sie beschreiben Bereiche, in denen das Elektron mit
hoher Wahrscheinlichkeit gefunden wird. Aus diesen Lösungen entsteht die
Vorstellung der Orbitale\index{Orbital}\index{Atomorbital}.

Damit unterscheidet sich das quantenmechanische Bild grundlegend vom bohrschen
Atommodell\index{bohrsches Atommodell}. Im bohrschen Modell bewegt sich das
Elektron noch auf bestimmten Bahnen um den Kern. In der Schrödinger-Theorie gibt
es solche Bahnen nicht mehr. Das Elektron wird durch einen Zustand beschrieben,
dessen räumliche Struktur durch die Wellenfunktion festgelegt ist.

Die Schrödingergleichung bildet damit den mathematischen Übergang vom
klassischen Elektronenbild zum modernen Atommodell. Sie erklärt sowohl die
diskreten Energien des Wasserstoffatoms als auch die räumliche Struktur der
Elektronenhülle\index{Elektronenhülle}.